\chapter{Python 解析の高速化}
\label{chap:performance}

\section{処理時間の基準測定}

高速化\index{こうそくか@高速化|textbf}を始める前に、どこが遅いのかを測る。読込みが遅いのか、探索が遅いのか、図の作成が遅いのかを区別しないまま最適化しても、効果は見えにくい。

ここで役立つのが \code{time.perf\_counter}\index{perf_counter@\code{time.perf\_counter}|textbf} である。処理の前後で時刻を取り、差を記録するだけでも、どの段階が重いかの見当が付く。大切なのは、全体時間だけでなく、読込み、前処理、探索、書き出し、作図のように段階を分けて測ることである。

\begin{lstlisting}[language=Python]
import time

t0 = time.perf_counter()
run_analysis()
t1 = time.perf_counter()
print(f"elapsed = {t1 - t0:.3f} s")
\end{lstlisting}

測定するときは、一度だけの結果を絶対視しないことも重要である。キャッシュや OS の状態で多少は揺れるため、条件をそろえ、複数回の傾向を見る方がよい。

\subsection{Jupyter の \code{\%timeit} による計測}

ノートブック環境で短い式や関数の速度を見たいときは、Jupyter の \code{\%timeit}\index{timeit@\code{\%timeit}|textbf} や \code{\%\%timeit} も便利である。複数回の実行を自動で繰り返し、平均的な実行時間を表示してくれるため、小さな差を見たいときに役立つ。

\begin{lstlisting}[language=Python]
%%timeit
result = array * 2
\end{lstlisting}

ただし、\code{\%timeit} はノートブックや IPython での簡便な測定であり、スクリプト全体の読込みや入出力を含めた時間を見るには向かない。小さな部品には \code{\%timeit}、解析全体には \code{perf\_counter} というように使い分けると整理しやすい。

\section{計算量の見積もり}
\index{けいさんりょう@計算量|textbf}
\index{あるごりずむ@アルゴリズム|textbf}

高速化では、実測だけでなく理屈の見積もりも役に立つ。入力サイズを $N$ としたとき、処理回数が $N$ に比例するのか、$N \log N$ 程度なのか、$N^2$ なのかで、大きなデータに対する伸び方がまるで違う。

厳密な漸近解析までできなくても、「データを 2 倍にしたら計算量は何倍くらいになるか」を考えるだけで十分に有益である。高速化はテクニック集ではなく、問題の構造を考える作業である。

変数名を短くしたり、関数を減らしたりしても、通常は実行時間にほとんど影響しない。一方、比較回数の削減、探索範囲の限定、配列演算への置換は、計算量または定数倍の処理時間を変え得る。

\subsection{オーダーとは何か}

ここでいう「オーダー」とは、入力サイズ $N$ を大きくしたときに処理回数がどのように増えるかを表す粗い記法である。通常は Big-O 記法\index{big-o@Big-O 記法|textbf}で書き、定数倍や低次の項を落として、最終的に支配的になる項だけを見る。

\[
3N + 20 \rightarrow O(N), \qquad
5N \log N + 100N \rightarrow O(N \log N), \qquad
N^2 + N \rightarrow O(N^2)
\]

よく出るオーダーには、たとえば次のようなものがある。

\begin{itemize}
\item $O(\log N)$: 範囲を半分ずつ絞る処理。二分探索など
\item $O(N)$: 1 回走査する処理。総和、平均、閾値判定、隣接差分など
\item $O(N \log N)$: ソートを含む処理。高速なソート、ソート後の探索など
\item $O(N^2)$: 全組合せ比較。全ペアの力ずく比較など
\end{itemize}

たとえば $N = 10^6$ なら、$\log_2 N$ はおよそ 20、$N \log_2 N$ はおよそ $2 \times 10^7$、$N^2$ は $10^{12}$ である。どちらも「増える」ことに変わりはないが、その増え方の差は圧倒的である。小さなテストデータでは目立たなくても、大きな入力では急に効いてくる。

ただし、オーダーは定数倍を完全に無視してよい、という意味ではない。同じ $O(N)$ でも 10 倍、100 倍の速度差が付くことは普通にある。ここでいう「ファクター」とは、この定数倍の差のことである。したがって、まずオーダーを疑い、その後で定数倍を詰める、という順序で考えると整理しやすい。

この中で、$O(\log N)$ や $O(N \log N)$ は初学者には直感を持ちにくい。そこで、まず二分探索で $O(\log N)$ を見てから、ソートで $O(N \log N)$ を見る。

\subsubsection{二分探索ではなぜ \texorpdfstring{$O(\log N)$}{O(log N)} になるか}

\(\log N\) は、「問題の大きさを毎回半分にできる」場面で現れる。たとえば昇順に並んだ列から目的の値を探す二分探索では、中央を見て、左半分か右半分のどちらかを捨てられる。

\begin{lstlisting}[language=Python]
def binary_search(values, target):
    left = 0
    right = len(values) - 1

    while left <= right:
        mid = (left + right) // 2
        if values[mid] == target:
            return mid
        if values[mid] < target:
            left = mid + 1
        else:
            right = mid - 1
    return -1
\end{lstlisting}

要素数が \(N\) 個あっても、1 回ごとに候補数が半分になるなら、必要な反復回数は \(N \to N/2 \to N/4 \to \cdots \to 1\) という形になる。これが \(\log_2 N\) 回程度である。したがって、二分探索は \(O(\log N)\) で動く。重要なのは、「1 個ずつ減らす」のではなく「毎回半分に減らす」ことが \(\log N\) を生む、という点である。

\subsubsection{ソートではなぜ \texorpdfstring{$O(N \log N)$}{O(N log N)} になるか}

ソートは、計算量の違いを理解するための基本例である。たとえば、挿入ソートやバブルソートのような初等的な方法は、一般に $O(N^2)$ で動く。一方、マージソートやヒープソートは $O(N \log N)$ であり、Python の \code{sorted()} も実用的な高速ソートを使っている。

ここで \(N \log N\) と書くなら、その意味も説明する必要がある。典型例であるマージソートでは、まず列を半分に分け、さらにその半分へ分ける、という分割を繰り返す。この分割の深さは、およそ \(\log_2 N\) 段である。一方、各段では、分かれた小列どうしを併合する過程で、全体として各要素を 1 回ずつ見るので \(O(N)\) かかる。

\[
\text{各段の仕事量 } O(N)
\quad \times \quad
\text{段数 } O(\log N)
\quad = \quad
O(N \log N)
\]

ここで重要なのは、「ソートは前処理にすぎない」と軽視しないことである。データを一度並べ替えるだけで、その後の探索や近傍判定がずっと安くなることがある。つまり、$O(N \log N)$ の前処理へ投資することで、後続の $O(N^2)$ を避けられるなら、全体として大幅に得をする。

データが 10 倍に増えたとき、$O(N^2)$ のソートはおおよそ 100 倍重くなるが、$O(N \log N)$ のソートは 10 倍強で済む。この差を理解しておくと、単にソートの名前を知るだけでなく、「なぜその方法を選ぶのか」を説明できるようになる。

\subsection{オーダーが改善する例: 重複判定}

問題をはっきりさせよう。長さ $N$ の整数列に、同じ値が 2 回以上現れるかを判定したいとする。

最も素朴な方法は、すべての組 $(i, j)$ を比べることである。

\begin{lstlisting}[language=Python]
def has_duplicate_slow(values):
    for i in range(len(values)):
        for j in range(i + 1, len(values)):
            if values[i] == values[j]:
                return True
    return False
\end{lstlisting}

この方法では、最悪の場合ほぼすべてのペアを比べるので $O(N^2)$ になる。データ量が 10 倍になると、比較回数はおおよそ 100 倍に増える。

しかし、いったんソートしてしまえば、同じ値どうしは必ず隣り合う。したがって、隣接要素だけを見れば重複を判定できる。

\begin{lstlisting}[language=Python]
def has_duplicate_fast(values):
    values_sorted = sorted(values)
    for i in range(len(values_sorted) - 1):
        if values_sorted[i] == values_sorted[i + 1]:
            return True
    return False
\end{lstlisting}

\code{sorted()} の部分は $O(N \log N)$、その後の 1 回走査は $O(N)$ なので、全体は $O(N \log N)$ である。元の $O(N^2)$ より明確に良い。ここで本質なのは、Python の書き方が少し変わったことではなく、「全比較」という問題の解き方をやめたことである。

Python の \code{set(values)} を使えば、重複判定は平均的に $O(N)$ で実行できる。ここでは、ソート後の走査によって処理量を減らす考え方を説明するため、ソートを用いた例を示している。

同じ考え方は探索でも現れる。未整列の列から目的の値を探す線形探索は $O(N)$ だが、整列済みの列なら二分探索で $O(\log N)$ にできる。したがって、「一度だけ探す」のか「何度も探す」のかによって、先にソートする価値があるかどうかも変わる。前処理に時間を使っても、その後の探索を何度も繰り返すなら、全体としては得になることが多い。

\subsection{オーダーは同じで定数倍だけ改善する例}

今度は、$N \times N$ 行列 $A$ が対称行列かどうか、すなわち $A_{ij} = A_{ji}$ が成り立つかを調べる問題を考える。

最も素朴な方法は、すべての $(i, j)$ について比較することである。

\begin{lstlisting}[language=Python]
def is_symmetric_slow(A):
    n = len(A)
    for i in range(n):
        for j in range(n):
            if A[i][j] != A[j][i]:
                return False
    return True
\end{lstlisting}

これは $n^2$ 回程度の比較を行うので $O(N^2)$ である。

しかし、対称性の判定では、上三角と下三角は同じ情報を二重に見ている。したがって、片側だけ見れば十分である。

\begin{lstlisting}[language=Python]
def is_symmetric_fast(A):
    n = len(A)
    for i in range(n):
        for j in range(i + 1, n):
            if A[i][j] != A[j][i]:
                return False
    return True
\end{lstlisting}

こちらもオーダーは同じく $O(N^2)$ だが、比較回数はおよそ $n(n-1)/2$ 回になり、素朴な実装のほぼ半分で済む。つまり、問題の大きさに対する増え方は変わっていないが、無駄な比較を消したことで定数倍が改善している。

この例は、「同じオーダーでも速くなる」ことを示している。図~\ref{fig:vectorization-benchmark} の隣接差分フィルタの比較も、基本的にはこの種類の改善である。比較回数そのものを 2 次から 1 次へ落としたのではなく、同じ種類の処理をより効率よく実装しているのである。

\section{ベクトル化・型・逐次処理による改善}

\subsection{NumPy と pandas による高速化}

アルゴリズムを見直した後は、NumPy や pandas へ処理を移すことで、Python レベルのループを減らせる。特に差分計算、ソート、列抽出、フィルタは、これらのライブラリが得意とする。図~\ref{fig:vectorization-benchmark} からも、ベクトル化\index{べくとるか@ベクトル化|textbf}が速度に与える効果を視覚的に確認できる。

ここで注意したいのは、NumPy や pandas を使えば自動的に何でも速くなるわけではないことである。根本的に無駄な比較をしているなら、配列演算へ移しても限界がある。まず計算量を減らし、そのうえで実装をベクトル化するという順序が自然である。

\subsection{メモリとコピーの管理}
\index{めもりかんり@メモリ管理|textbf}
\index{びゅー@ビュー|textbf}
\index{こぴー@コピー}

速度だけを見ていると、メモリ使用量が大きくなりすぎて逆に遅くなることがある。巨大な中間配列を何度も作る実装では、CPU よりもメモリ転送が支配的になる場合がある。

そのため、高速化では「何回ループしたか」だけでなく、「何回コピーしたか」も見る必要がある。不要なコピーを減らし、必要な列や範囲だけを扱う設計は、速度と安定性の両方に効く。

NumPy では、スライスが元配列とメモリを共有するビューになる一方、条件抽出は新しい配列を作る。次のコードで両者を確認できる。

\begin{lstlisting}[language=Python]
import numpy as np

values = np.arange(1_000_000, dtype=np.float64)
view = values[100:200]
copied = values[values % 2 == 0]

print(values.nbytes)
print(view.nbytes, np.shares_memory(values, view))
print(copied.nbytes, np.shares_memory(values, copied))
\end{lstlisting}

実行結果は次の通りである。

\begin{lstlisting}[numbers=none, xleftmargin=2\zw]
8000000
800 True
4000000 False
\end{lstlisting}

\code{nbytes} は配列要素が占めるバイト数を表す。\code{np.shares\_memory(a, b)} は 2 配列が同じメモリ領域を共有するかを調べる。この例では、\code{view} 自体が参照する要素は 800 バイト分だが、元配列のメモリを共有するため、\code{view} の要素を書き換えると対応する \code{values} の要素も変わる。\code{copied} は 50 万要素の独立した配列であり、新たに 4 MB を使う。なお、\code{nbytes} は Python オブジェクト自体の管理領域を含まないため、プロセス全体のメモリ使用量と完全には一致しない。

\subsection{分割読込み、並列化、ストリーミング}
\index{ぶんかつよみこみ@分割読込み|textbf}
\index{へいれつか@並列化|textbf}
\index{すとりーみんぐ@ストリーミング|textbf}

大規模解析では、全部を一度にメモリへ載せない設計も重要である。ファイル単位で処理し、必要な情報だけを持ち回る。さらに、独立な単位に分割できるなら並列化も検討できる。

ストリーミング処理では、「今必要な範囲だけを覚える」という考え方が中心になる。全履歴が不要なら保持しない、集計済みの量は逐次更新する、といった設計が有効である。メモリ節約と速度向上が同時に得られることも多い。

例えば、1 行に 1 個の測定値が書かれた \code{values.txt} の平均は、全行をリストへ保存せずに求められる。

\begin{lstlisting}[numbers=none, xleftmargin=2\zw, title={values.txt}]
1.5
2.0
4.0
2.5
\end{lstlisting}

\begin{lstlisting}[language=Python]
def streaming_mean(path):
    count = 0
    total = 0.0

    with open(path, "r", encoding="utf-8") as f:
        for line_number, line in enumerate(f, start=1):
            text = line.strip()
            if text == "":
                continue
            try:
                value = float(text)
            except ValueError as exc:
                raise ValueError(f"line {line_number}: {text!r}") from exc
            count += 1
            total += value

    if count == 0:
        raise ValueError("input contains no numeric values")
    return total / count


print(streaming_mean("values.txt"))
\end{lstlisting}

出力は \code{2.5} となる。引数 \code{path} は入力ファイルの場所であり、\code{enumerate(..., start=1)} はエラー箇所をファイルと同じ 1 始まりの行番号で報告するために使う。この関数が保持するのは現在行、件数、合計だけなので、入力行数が増えても集計用メモリはほぼ一定である。ただし、中央値や全体のソートのように全データの順序が必要な統計量は、この単純な方法だけでは正確に求められない。

\section{プロファイリングと低レベル実装}

\subsection{プロファイラによるボトルネックの特定}
\index{ぼとるねっく@ボトルネック|textbf}

本格的に遅い理由を調べるなら、プロファイラ\index{ぷろふぁいら@プロファイラ|textbf}も役立つ。例えば \code{cProfile}\index{cprofile@\code{cProfile}|textbf} を使うと、どの関数に時間が使われているかを確認できる。

\begin{lstlisting}[language=Python]
import cProfile

cProfile.run("run_analysis()", sort="cumtime")
\end{lstlisting}

ただし、プロファイラの数字も文脈なしには読めない。回数が多いから遅いのか、1 回が重いのか、入出力を待っているのかを区別する必要がある。測定値とコード構造を一緒に見ることが重要である。

\begin{lstlisting}[numbers=none, xleftmargin=2\zw]
ncalls  tottime  cumtime  function
100000   1.25     1.40    check_pair
    10   0.02     2.10    read_all_files
\end{lstlisting}

この結果例では、\code{check\_pair} は 1 回は軽いが回数が多く、\code{read\_all\_files} は内部の処理込みで累積時間が大きいと読める。どの列が何を意味するかを理解すると、改善方針が立てやすくなる。

\subsection{C や C++ へ書き換えるという選択肢}

Python は開発速度が高いが、計算コアのすべてを Python で書く必要はない。アルゴリズムを十分に整理したうえで、どうしても重い部分だけを C や C++ へ移す、というのは自然な選択肢である。

ただし、言語を変える前に Python 側で何が遅いかを把握しておく必要がある。ボトルネックが入出力や高水準の設計にあるなら、C++ 化しても大きな改善にならない。低レベル言語への移行は、最後の切り札として考える方がよい。

\begin{lstlisting}[language=C++]
for (std::size_t i = 0; i < n; ++i) {
    for (std::size_t j = i + 1; j < n; ++j) {
        if (std::abs(times[j] - times[i]) <= window) {
            ++count;
        }
    }
}
\end{lstlisting}

この書き換えは、定数倍の改善をもたらすことがある。しかし、比較回数そのものが多すぎるなら本質的な改善ではない。アルゴリズムと実装言語を分けて考える姿勢が必要である。

\subsection{Python からの C/C++ コード呼び出し}
\index{cython@Cython|textbf}
\index{pybind11@pybind11|textbf}

方法はいくつかある。たとえば \code{ctypes} で共有ライブラリを呼ぶ方法、Python C 拡張を書く方法、pybind11 を使って C++ を Python から自然に使える形へ包む方法などがある。さらに、Python と C の中間に位置する Cython を使い、Python に近い構文のまま重いループだけをコンパイルする方法も、この節に含まれる重要な選択肢である。

\begin{lstlisting}[language=Python]
cdef int i
cdef double total = 0.0

for i in range(n):
    total += values[i]
\end{lstlisting}

ここで \code{cdef} は Cython 特有の宣言であり、変数を Python オブジェクトではなく C の整数や浮動小数点数として扱わせるために使う。これにより、ループ内部の型判定やボックス化のコストを減らせる。

どの方法を選ぶかは、必要な速度、保守性、ビルドの難しさで決まる。目安としては次のように考えるとよい。

\begin{itemize}
\item \code{ctypes}: 既にある C の共有ライブラリを手早く呼びたいときに向く。設定は比較的軽いが、引数や戻り値の型を手で合わせる必要がある
\item Python C 拡張: Python 本体の C API を直接使う最も低レベルな方法で、自由度は高いが記述量も多い
\item \code{pybind11}: C++ の関数やクラスを Python から自然に使える形で公開しやすい。C++ 側の実装をすでに持っている場合に相性がよい
\item Cython: Python に近い構文を保ちながら、重いループや数値計算だけへ型を付けてコンパイルしたいときに向く
\end{itemize}

短い試作なら \code{ctypes}、Python らしい使い勝手まで整えたいなら pybind11、既存の Python 風コードを保ちながら内側だけ速くしたいなら Cython、といった選択が考えられる。重要なのは、インタフェースを狭く保ち、重い部分だけを外へ出すことである。pybind11 の最小構成は付録~\ref{app:pybind11} を参照してほしい。

\begin{lstlisting}[language=Python]
result = fast_module.count_pairs(times, window)
\end{lstlisting}

Python 側の呼び出しが単純であれば、内部実装を後から差し替えやすい。まず Python で正しい API を決め、その内部だけを高速化するという順が自然である。なお、Cython も強力だが、アルゴリズムそのものが無駄な場合には決定打にならない。この点は C++ 化と同じであり、まず Python 側で何が遅いのか、どのループが本当に支配的なのかを見極めた後に使う方が効果的である。

\subsection{外部実装を参考にするときの見方}

高速化を学ぶときは、外部の C や C++ の実装例を読むことが参考になる。ただし、そこで大切なのは固有名詞を覚えることではない。まず見るべきなのは、どの問題を解いているのか、どのアルゴリズムを選んでいるのか、どの部分だけを低レベル言語へ移しているのか、という点である。

高速化の章では、このような例を参照しながら、次の二つを常に区別して考える。

\begin{itemize}
\item アルゴリズムを変えたから速くなったのか
\item 実装言語を変えたから速くなったのか
\end{itemize}

多くの場合、最初に効くのは前者である。

したがって、高速化のレポートを書くときは、「NumPy にした」「C++ にした」で終わらせない方がよい。何が変わった結果として速くなったのかを、計算量、比較回数、メモリ配置、ループの位置といった観点で説明するべきである。

\section{高速化前後の比較と測定条件の報告}

高速化の結果は、本文中で比較できる形に整理すると読みやすい。例えば、次は報告形式を示すための仮想的な数値例である。

\begin{lstlisting}[numbers=none, xleftmargin=2\zw]
slow version : 182.4 s
vectorized   :  12.7 s
streaming    :   4.8 s
\end{lstlisting}

実際の報告では、この形式へ自分で測定した値を入れる。この表現だけでは理由が分からないので、続けて「比較回数を減らした」「Python ループを配列演算へ移した」「全件保持をやめた」といった説明を添える。数値と理由の両方がそろって初めて報告になる。

特に、高速化の報告では「$O(N^2)$ から $O(N \log N)$ や $O(N)$ へ落とした改善」と、「同じ $O(N)$ や $O(N^2)$ のまま NumPy, Cython, C++ などで定数倍を詰めた改善」を分けて書くとよい。前者は問題の解き方を変えた改善であり、後者は実装の効率を上げた改善である。この二つを混同しない方が、読者にも自分にも何が効いたのかが伝わりやすい。

\begin{lstlisting}[language=Python]
import numpy as np

values = np.array(data, dtype=np.float64)
diffs = np.diff(values)
count = np.count_nonzero(diffs > 0.25)
\end{lstlisting}

\begin{figure}[tbp]
\centering
\includegraphics[width=0.82\linewidth]{vectorization_benchmark.pdf}
\caption{実行時間の比較方法を示す模式例。値は説明用であり、特定の計算機で得たベンチマーク結果ではない。縦軸には対数目盛を用いている。}
\label{fig:vectorization-benchmark}
\end{figure}

図~\ref{fig:vectorization-benchmark} は実行時間の大小を比較するための図であり、結果の正しさまでは保証しない。高速化の前後で同じ \code{count} が得られることをテストし、その上で入力サイズ、反復回数、実行環境、代表値をそろえて時間を比較する必要がある。

\section{処理時間・メモリ・プロファイルの測定演習}

\begin{enumerate}
\item \textbf{結果を保持したベクトル化}。
\verb|values = np.arange(1_000_000, dtype=np.float64)| について、各要素を二乗してから総和を求める処理を Python の \code{for} 文と NumPy の列演算で実装せよ。\code{np.isclose()} で結果が一致することを先に確認し、その後で \code{time.perf\_counter()} を使って各方法を 5 回ずつ測定せよ。最小値または中央値を報告し、1 回だけの値を比較しない理由を説明せよ。

\item \textbf{入力サイズと計算量の関係}。
本文の \code{has\_duplicate\_slow()} と \code{has\_duplicate\_fast()} を使い、重複を含まない入力について \(N=1000, 2000, 4000\) の実行時間を測定せよ。\(N\) を 2 倍にしたときの時間比を求め、\(O(N^2)\) と \(O(N\log N)\) の予想と比較せよ。計算機や実行時の状態によって時間は変わるため、絶対時間ではなく比と傾向を説明せよ。

\item \textbf{ビューとコピーの影響}。
本文の \code{values}、\code{view}、\code{copied} を作り、\code{view[0] = -1} と \code{copied[0] = -2} を実行した後で \code{values[100]} と \code{values[0]} を表示せよ。前者だけが変化する理由を、\code{np.shares\_memory()} の結果と対応させて説明せよ。さらに \code{float32} と \code{float64} で同じ要素数の \code{nbytes} を比較せよ。

\item \textbf{逐次集計の検証}。
1 から 100000 までを 1 行ずつ書いたファイルを作り、本文の \code{streaming\_mean()} で平均を求めよ。期待値は \code{50000.5} である。同じファイルを全件リストへ読み込む実装も作り、結果の一致を確認せよ。両者が保持するデータ量を、入力行数 \(N\) に対するオーダーで説明せよ。

\item \textbf{プロファイル結果に基づく改善対象の選定}。
自分の解析スクリプト、または 1 問目の 2 実装を関数へ分けたスクリプトを \code{python -m cProfile -s cumtime script.py} で測定せよ。出力の \code{ncalls}、\code{tottime}、\code{cumtime} の意味を説明し、累積時間の上位 3 関数を記録せよ。そのうち改善対象に選ぶ関数と、選ばない関数の理由をそれぞれ述べよ。
\end{enumerate}

性能改善は、速いコードを書くことだけを目的としない。高速化の前後で同じ結果が得られることを検証し、入力条件と測定方法を記録して初めて、改善を再現できる。次章では、同時観測イベントの探索を題材として、読込み、検証、可視化、フィット、性能測定を一つの解析へ統合する。
